% --- Archon physics formalization source begin ---
% archon:physics
% archon:covers PhyXMiniProblems/problem_phyx_mini_0817.lean
% archon:source-report reports/phyx_mini/problem_phyx_mini_0817.source.json

\chapter{Physics problem phyx\_mini\_0817}
\label{ch:PhyXMiniProblems_problem_phyx_mini_0817}

\paragraph{Problem source.}
\#\# Physical scenario

James (mass 90.0 kg) and Ramon (mass 60.0 kg) are 20.0 m apart on a frozen pond. Midway between them is a mug of their favorite beverage as shown in figure. They pull on the ends of a light rope stretched between them.

\#\# Question

When James has moved 6.0 m toward the mug, how far has Ramon moved?

\#\# Answer choices

- A: A: 6.0m

- B: B: 9.0m

- C: C: 4.0m

- D: D: 3.0m

\#\# Figure caption (auxiliary; use the image as primary evidence)

The image is a diagram with two stick figures on either side, labeled "James" on the left and "Ramon" on the right. 

- **James** is on the left side, positioned at -10.0 meters on the horizontal axis and has a mass of 90.0 kg.
- **Ramon** is on the right side, positioned at +10.0 meters and has a mass of 60.0 kg.
  
The axis is marked with various points:
- \textbackslash{}( x\_\{cm\} \textbackslash{}) is marked along the axis to indicate the center of mass, although its exact position isn't specified in the image.
- A point labeled "P" is positioned at 0 on the axis.

The stick figures are each holding one end of a line that spans the distance between them, indicating the presence of some kind of connection or relationship between the two points.

\#\# Dataset metadata

Momentum and Collisions; Spatial Relation Reasoning, Physical Model Grounding Reasoning

\paragraph{Recorded answer/context.}
B: B: 9.0m

\paragraph{Figure/image path.}
/root/proposal\_for\_physic/science-mango/phyx\_mini\_run/phyx\_data/test\_image/817.png

\paragraph{Formalization target.}
create a compiling Lean file with sorry bodies at `PhyXMiniProblems/problem\_phyx\_mini\_0817.lean`. The Lean declarations must preserve the physical quantities, dimensions or dimensional roles, figure labels, governing-law hypotheses, and final relation expressed by this problem.
Use Mathlib/Physlib names found through LeanExplore where available. If a physics API is missing, introduce faithful local abstractions rather than scalar placeholder aliases.

\begin{theorem}[Physics formalization target]
\label{thm:physics:phyx_mini_0817:target}
\lean{PhyXMiniProblems.ProblemPhyXMini0817.ramon_distanceMoved_eq_nine_meters_and_choiceB}
\uses{def:physics:phyx-mini-0817:phyxminiproblems-problemphyxmini0817-frozenpondropesetup, def:physics:phyx-mini-0817:phyxminiproblems-problemphyxmini0817-distancemovedinmeters, def:physics:phyx-mini-0817:phyxminiproblems-problemphyxmini0817-matchesproblemandprimaryfigure, def:physics:phyx-mini-0817:phyxminiproblems-problemphyxmini0817-whenjameshasmovedsixmeterstowardmug, def:physics:phyx-mini-0817:phyxminiproblems-problemphyxmini0817-satisfiesisolatedhorizontalcenterofmasslaw, def:physics:phyx-mini-0817:phyxminiproblems-problemphyxmini0817-answerchoice, def:physics:phyx-mini-0817:phyxminiproblems-problemphyxmini0817-answerchoice-distanceinmeters}
The assigned autoformalize agent should translate this physics problem into Lean declarations in the covered file, with theorem and lemma proof bodies written as `by sorry`.
\end{theorem}
\begin{proof}
This is an autoformalization task, not a proof task. Produce statements that can later be proved by the physics prover without weakening the physical model.
\end{proof}

% --- Archon declaration topology begin ---
\section{Declaration topology}
\begin{definition}[Mass Quantity]
  \label{def:physics:phyx-mini-0817:phyxminiproblems-problemphyxmini0817-massquantity}
  \lean{PhyXMiniProblems.ProblemPhyXMini0817.MassQuantity}
  A nonnegative physical mass, independent of the chosen unit system
\end{definition}
\begin{proof}
  This is the declared model definition of Mass Quantity.
\end{proof}
\begin{definition}[Axial Position]
  \label{def:physics:phyx-mini-0817:phyxminiproblems-problemphyxmini0817-axialposition}
  \lean{PhyXMiniProblems.ProblemPhyXMini0817.AxialPosition}
  A signed physical coordinate on the horizontal axis in the figure
\end{definition}
\begin{proof}
  This is the declared model definition of Axial Position.
\end{proof}
\begin{definition}[Length Quantity]
  \label{def:physics:phyx-mini-0817:phyxminiproblems-problemphyxmini0817-lengthquantity}
  \lean{PhyXMiniProblems.ProblemPhyXMini0817.LengthQuantity}
  A nonnegative physical length
\end{definition}
\begin{proof}
  This is the declared model definition of Length Quantity.
\end{proof}
\begin{definition}[mass In Kilograms]
  \label{def:physics:phyx-mini-0817:phyxminiproblems-problemphyxmini0817-massinkilograms}
  \lean{PhyXMiniProblems.ProblemPhyXMini0817.massInKilograms}
  \uses{def:physics:phyx-mini-0817:phyxminiproblems-problemphyxmini0817-massquantity}
  Kilogram readout of a physical mass
\end{definition}
\begin{proof}
  This is the declared model definition of mass In Kilograms.
\end{proof}
\begin{definition}[position In Meters]
  \label{def:physics:phyx-mini-0817:phyxminiproblems-problemphyxmini0817-positioninmeters}
  \lean{PhyXMiniProblems.ProblemPhyXMini0817.positionInMeters}
  \uses{def:physics:phyx-mini-0817:phyxminiproblems-problemphyxmini0817-axialposition}
  Metre readout of a signed horizontal position
\end{definition}
\begin{proof}
  This is the declared model definition of position In Meters.
\end{proof}
\begin{definition}[length In Meters]
  \label{def:physics:phyx-mini-0817:phyxminiproblems-problemphyxmini0817-lengthinmeters}
  \lean{PhyXMiniProblems.ProblemPhyXMini0817.lengthInMeters}
  \uses{def:physics:phyx-mini-0817:phyxminiproblems-problemphyxmini0817-lengthquantity}
  Metre readout of a nonnegative physical length
\end{definition}
\begin{proof}
  This is the declared model definition of length In Meters.
\end{proof}
\begin{definition}[Person]
  \label{def:physics:phyx-mini-0817:phyxminiproblems-problemphyxmini0817-person}
  \lean{PhyXMiniProblems.ProblemPhyXMini0817.Person}
  The two people named in the problem and primary figure
\end{definition}
\begin{proof}
  This is the declared model definition of Person.
\end{proof}
\begin{definition}[Observation Time]
  \label{def:physics:phyx-mini-0817:phyxminiproblems-problemphyxmini0817-observationtime}
  \lean{PhyXMiniProblems.ProblemPhyXMini0817.ObservationTime}
  The two instants compared in the question
\end{definition}
\begin{proof}
  This is the declared model definition of Observation Time.
\end{proof}
\begin{definition}[Axis Marker]
  \label{def:physics:phyx-mini-0817:phyxminiproblems-problemphyxmini0817-axismarker}
  \lean{PhyXMiniProblems.ProblemPhyXMini0817.AxisMarker}
  Labels printed on the horizontal axis of the primary figure
\end{definition}
\begin{proof}
  This is the declared model definition of Axis Marker.
\end{proof}
\begin{definition}[Rope Mass Model]
  \label{def:physics:phyx-mini-0817:phyxminiproblems-problemphyxmini0817-ropemassmodel}
  \lean{PhyXMiniProblems.ProblemPhyXMini0817.RopeMassModel}
  The light-rope idealization used in the mechanics model
\end{definition}
\begin{proof}
  This is the declared model definition of Rope Mass Model.
\end{proof}
\begin{definition}[Rope Configuration]
  \label{def:physics:phyx-mini-0817:phyxminiproblems-problemphyxmini0817-ropeconfiguration}
  \lean{PhyXMiniProblems.ProblemPhyXMini0817.RopeConfiguration}
  The rope configuration described and drawn in the problem
\end{definition}
\begin{proof}
  This is the declared model definition of Rope Configuration.
\end{proof}
\begin{definition}[Rope Interaction]
  \label{def:physics:phyx-mini-0817:phyxminiproblems-problemphyxmini0817-ropeinteraction}
  \lean{PhyXMiniProblems.ProblemPhyXMini0817.RopeInteraction}
  The interaction described by “they pull on the ends of a light rope.”
\end{definition}
\begin{proof}
  This is the declared model definition of Rope Interaction.
\end{proof}
\begin{definition}[Pond Surface Model]
  \label{def:physics:phyx-mini-0817:phyxminiproblems-problemphyxmini0817-pondsurfacemodel}
  \lean{PhyXMiniProblems.ProblemPhyXMini0817.PondSurfaceModel}
  The horizontal idealization of motion on the frozen pond
\end{definition}
\begin{proof}
  This is the declared model definition of Pond Surface Model.
\end{proof}
\begin{definition}[Initial Motion Model]
  \label{def:physics:phyx-mini-0817:phyxminiproblems-problemphyxmini0817-initialmotionmodel}
  \lean{PhyXMiniProblems.ProblemPhyXMini0817.InitialMotionModel}
  The initial motion state implicit before James and Ramon begin pulling
\end{definition}
\begin{proof}
  This is the declared model definition of Initial Motion Model.
\end{proof}
\begin{definition}[Frozen Pond Rope Setup]
  \label{def:physics:phyx-mini-0817:phyxminiproblems-problemphyxmini0817-frozenpondropesetup}
  \lean{PhyXMiniProblems.ProblemPhyXMini0817.FrozenPondRopeSetup}
  \uses{def:physics:phyx-mini-0817:phyxminiproblems-problemphyxmini0817-massquantity, def:physics:phyx-mini-0817:phyxminiproblems-problemphyxmini0817-axialposition, def:physics:phyx-mini-0817:phyxminiproblems-problemphyxmini0817-lengthquantity, def:physics:phyx-mini-0817:phyxminiproblems-problemphyxmini0817-person, def:physics:phyx-mini-0817:phyxminiproblems-problemphyxmini0817-observationtime, def:physics:phyx-mini-0817:phyxminiproblems-problemphyxmini0817-axismarker, def:physics:phyx-mini-0817:phyxminiproblems-problemphyxmini0817-ropemassmodel, def:physics:phyx-mini-0817:phyxminiproblems-problemphyxmini0817-ropeconfiguration, def:physics:phyx-mini-0817:phyxminiproblems-problemphyxmini0817-ropeinteraction, def:physics:phyx-mini-0817:phyxminiproblems-problemphyxmini0817-pondsurfacemodel, def:physics:phyx-mini-0817:phyxminiproblems-problemphyxmini0817-initialmotionmodel}
  All independent physical quantities and qualitative labels in the pictured two-person system. In particular, Ramon's later position is an independent field and is not assigned its requested value here
\end{definition}
\begin{proof}
  This is the declared model definition of Frozen Pond Rope Setup.
\end{proof}
\begin{definition}[two Person Center Of Mass In Meters]
  \label{def:physics:phyx-mini-0817:phyxminiproblems-problemphyxmini0817-twopersoncenterofmassinmeters}
  \lean{PhyXMiniProblems.ProblemPhyXMini0817.twoPersonCenterOfMassInMeters}
  \uses{def:physics:phyx-mini-0817:phyxminiproblems-problemphyxmini0817-massinkilograms, def:physics:phyx-mini-0817:phyxminiproblems-problemphyxmini0817-positioninmeters, def:physics:phyx-mini-0817:phyxminiproblems-problemphyxmini0817-person, def:physics:phyx-mini-0817:phyxminiproblems-problemphyxmini0817-observationtime, def:physics:phyx-mini-0817:phyxminiproblems-problemphyxmini0817-frozenpondropesetup}
  The center-of-mass coordinate of James and Ramon at one observation time, using their kilogram and metre readouts. This is Mathlib's finite weighted center of mass specialized to the two named people
\end{definition}
\begin{proof}
  This is the declared model definition of two Person Center Of Mass In Meters.
\end{proof}
\begin{definition}[distance Moved In Meters]
  \label{def:physics:phyx-mini-0817:phyxminiproblems-problemphyxmini0817-distancemovedinmeters}
  \lean{PhyXMiniProblems.ProblemPhyXMini0817.distanceMovedInMeters}
  \uses{def:physics:phyx-mini-0817:phyxminiproblems-problemphyxmini0817-positioninmeters, def:physics:phyx-mini-0817:phyxminiproblems-problemphyxmini0817-person, def:physics:phyx-mini-0817:phyxminiproblems-problemphyxmini0817-frozenpondropesetup}
  The ordinary nonnegative distance a named person travels, in metres
\end{definition}
\begin{proof}
  This is the declared model definition of distance Moved In Meters.
\end{proof}
\begin{definition}[Matches Problem And Primary Figure]
  \label{def:physics:phyx-mini-0817:phyxminiproblems-problemphyxmini0817-matchesproblemandprimaryfigure}
  \lean{PhyXMiniProblems.ProblemPhyXMini0817.MatchesProblemAndPrimaryFigure}
  \uses{def:physics:phyx-mini-0817:phyxminiproblems-problemphyxmini0817-massinkilograms, def:physics:phyx-mini-0817:phyxminiproblems-problemphyxmini0817-positioninmeters, def:physics:phyx-mini-0817:phyxminiproblems-problemphyxmini0817-lengthinmeters, def:physics:phyx-mini-0817:phyxminiproblems-problemphyxmini0817-frozenpondropesetup, def:physics:phyx-mini-0817:phyxminiproblems-problemphyxmini0817-twopersoncenterofmassinmeters}
  Numerical and qualitative readouts supplied by the prose and primary image. The midpoint equation records that the mug is midway between the initial positions. The x\_cm marker is identified with the initial center of mass without assigning it an extra numerical value
\end{definition}
\begin{proof}
  This is the declared model definition of Matches Problem And Primary Figure.
\end{proof}
\begin{definition}[When James Has Moved Six Meters Toward Mug]
  \label{def:physics:phyx-mini-0817:phyxminiproblems-problemphyxmini0817-whenjameshasmovedsixmeterstowardmug}
  \lean{PhyXMiniProblems.ProblemPhyXMini0817.WhenJamesHasMovedSixMetersTowardMug}
  \uses{def:physics:phyx-mini-0817:phyxminiproblems-problemphyxmini0817-positioninmeters, def:physics:phyx-mini-0817:phyxminiproblems-problemphyxmini0817-frozenpondropesetup}
  The observation specified in the question: James has moved six metres in the positive-axis direction from his pictured location toward P, without passing the mug. This contains no condition on Ramon's later position
\end{definition}
\begin{proof}
  This is the declared model definition of When James Has Moved Six Meters Toward Mug.
\end{proof}
\begin{definition}[Satisfies Isolated Horizontal Center Of Mass Law]
  \label{def:physics:phyx-mini-0817:phyxminiproblems-problemphyxmini0817-satisfiesisolatedhorizontalcenterofmasslaw}
  \lean{PhyXMiniProblems.ProblemPhyXMini0817.SatisfiesIsolatedHorizontalCenterOfMassLaw}
  \uses{def:physics:phyx-mini-0817:phyxminiproblems-problemphyxmini0817-frozenpondropesetup, def:physics:phyx-mini-0817:phyxminiproblems-problemphyxmini0817-twopersoncenterofmassinmeters}
  For the light-rope system on the frictionless horizontal pond, all horizontal pulls are internal. Starting from rest gives zero initial total horizontal momentum, so the two-person center of mass is unchanged between the two observation times
\end{definition}
\begin{proof}
  This is the declared model definition of Satisfies Isolated Horizontal Center Of Mass Law.
\end{proof}
\begin{theorem}[mass Weighted Displacements balance]
  \label{thm:physics:phyx-mini-0817:phyxminiproblems-problemphyxmini0817-massweighteddisplacements-balance}
  \lean{PhyXMiniProblems.ProblemPhyXMini0817.massWeightedDisplacements_balance}
  \uses{def:physics:phyx-mini-0817:phyxminiproblems-problemphyxmini0817-massinkilograms, def:physics:phyx-mini-0817:phyxminiproblems-problemphyxmini0817-positioninmeters, def:physics:phyx-mini-0817:phyxminiproblems-problemphyxmini0817-frozenpondropesetup, def:physics:phyx-mini-0817:phyxminiproblems-problemphyxmini0817-matchesproblemandprimaryfigure, def:physics:phyx-mini-0817:phyxminiproblems-problemphyxmini0817-satisfiesisolatedhorizontalcenterofmasslaw}
  Center-of-mass conservation gives the mass-weighted balance of the two signed axial displacements. This intermediate relation is derived rather than included in the governing-law premise
\end{theorem}
\begin{proof}
  The conclusion follows from the governing relations and figure readouts named in the statement of mass Weighted Displacements balance.
\end{proof}
\begin{definition}[Answer Choice]
  \label{def:physics:phyx-mini-0817:phyxminiproblems-problemphyxmini0817-answerchoice}
  \lean{PhyXMiniProblems.ProblemPhyXMini0817.AnswerChoice}
  Labels of the four multiple-choice answers
\end{definition}
\begin{proof}
  This is the declared model definition of Answer Choice.
\end{proof}
\begin{definition}[distance In Meters]
  \label{def:physics:phyx-mini-0817:phyxminiproblems-problemphyxmini0817-answerchoice-distanceinmeters}
  \lean{PhyXMiniProblems.ProblemPhyXMini0817.AnswerChoice.distanceInMeters}
  \uses{def:physics:phyx-mini-0817:phyxminiproblems-problemphyxmini0817-answerchoice}
  Distance in metres printed beside each displayed answer label
\end{definition}
\begin{proof}
  This is the declared model definition of distance In Meters.
\end{proof}
% --- Archon declaration topology end ---
% --- Archon physics formalization source end ---
